% Generated draft: requires mathematical and source review.
\documentclass[11pt]{article}
\usepackage[margin=25mm]{geometry}
\usepackage{fontspec}
\setmainfont{Latin Modern Roman}
\usepackage{amsmath,amssymb,mathtools,unicode-math}
\setmathfont{Latin Modern Math}
\usepackage{enumitem,needspace,xurl,hyperref}
\setlength{\parindent}{0pt}
\setlength{\parskip}{5pt}
\setlength{\emergencystretch}{3em}
\newcommand{\N}{\mathbb{N}}
\newcommand{\Z}{\mathbb{Z}}
\newcommand{\Q}{\mathbb{Q}}
\newcommand{\R}{\mathbb{R}}
\newcommand{\C}{\mathbb{C}}
\newcommand{\F}{\mathbb{F}}
\newcommand{\A}{\mathbb{A}}
\newcommand{\PP}{\mathbb{P}}
\newcommand{\E}{\mathbb{E}}
\newcommand{\eps}{\varepsilon}
\newcommand{\dd}{\,\mathrm{d}}
\begin{document}
\section*{UnsolvedMath ID: 800406}
\textbf{Source problem code:} catalogue entry
\subsection{Definitions and mathematical statement}
Let \(f:\PP^N\dashrightarrow\PP^N\) be a dominant rational map defined over \(\overline{\Q}\). Write \(I_f\) for its indeterminacy locus and define
\[
\PP^N(\overline{\Q})_f=\{P\in\PP^N(\overline{\Q}): f^n(P)\text{ is defined for every }n\geq 0\},
\qquad
\mathcal O_f(P)=\{f^n(P):n\geq0\}.
\]
In positive dimension, the usual first dynamical degree is
\[
\delta_f=\lim_{m\to\infty}\bigl(\deg(f^m)\bigr)^{1/m},
\]
where the degree of a rational map is the common degree of homogeneous coordinate polynomials after cancellation of common factors. The excerpt does not explicitly specify the range of \(N\); this definition uses the usual setting \(N\geq1\).

The supplied statement uses a height \(h\) without defining it. A conventional interpretation is the absolute logarithmic Weil height on \(\PP^N(\overline{\Q})\): for homogeneous coordinates in a number field \(K\),
\[
h([x_0:\cdots:x_N])=
\frac{1}{[K:\Q]}\sum_{v\in M_K}[K_v:\Q_v]\log\max_{0\leq i\leq N}|x_i|_v,
\]
with the usual normalized absolute values extending those of \(\Q\). This interpretation requires confirmation from the source; the excerpt does not specify whether a modified height such as \(\max\{1,h\}\) is intended.

\textbf{Question (supplied source, Question 29).} Suppose \(P\in\PP^N(\overline{\Q})_f\) satisfies
\[
\overline{\mathcal O_f(P)}^{\mathrm{Zar}}=\PP^N.
\]
Do there exist integers \(k_f\in\Z\) and \(\ell_f\in\{0,1,\ldots,N\}\) such that
\[
h(f^n(P))\asymp \delta_f^n n^{\ell_f}(\log n)^{k_f}
\qquad(n\to\infty)?
\]
Explicitly, the requested comparison means that there are constants \(c,C>0\) and an integer \(n_0\geq2\), independent of \(n\), such that for every integer \(n\geq n_0\),
\[
c\,\delta_f^n n^{\ell_f}(\log n)^{k_f}
\leq h(f^n(P))
\leq C\,\delta_f^n n^{\ell_f}(\log n)^{k_f}.
\]
The source expressly allows the comparison constants to depend on \(f\) and \(P\). Further, when \(\delta_f>1\), can the comparison always be obtained with \(k_f=0\)?

The displayed subscripts suggest exponents attached to \(f\), but the supplied wording first fixes both \(f\) and \(P\). It does not explicitly resolve whether one pair of exponents must work for every Zariski-dense orbit of a fixed map, or whether the exponents may also depend on \(P\). This distinction remains unresolved here.
\subsection{Short English statement}
For a dominant rational self-map of projective space over \(\overline{\Q}\), ask whether the height along every everywhere-defined, Zariski-dense orbit is comparable to exponential growth given by the dynamical degree, multiplied by an integer power of \(n\) between \(0\) and \(N\) and an integer power of \(\log n\). If the dynamical degree exceeds \(1\), ask whether the logarithmic factor can always be omitted. The height convention and whether the exponents must be uniform over dense orbits of a fixed map require clarification.
\subsection{Sources}
\begin{itemize}
\item Supplied local UnsolvedMath registry record: permanent ID \(800406\), exact source problem code \texttt{catalogue entry}, title ``Open Problems in Arithmetic Dynamics --- Question 29.'' The supplied registry labels the problem \texttt{open}; this is a source label, not an independently verified mathematical or literature assessment.
\item Benedetto, DeMarco, Ingram, Jones et al., \emph{Current Trends and Open Problems in Arithmetic Dynamics} (2018), Question 29. Supplied source link: \url{https://arxiv.org/abs/1806.04980}. The registry attributes the proposal to these authors and gives the year 2018.
\item Supplied provenance: source accessed on 2026-10-16; extraction method \texttt{arxiv\_latex}. These are supplied provenance claims, not a claim of source access or verification for this draft. The registry record was created and updated on 2026-10-16 and records upstream revision \texttt{d371ee917ac9647fbb06cd6fc529deec6f3dda82}.
\end{itemize}
\end{document}
